
\section{Definitions}
We begin with some definitions, following the notation in
\cite{KNRW18}.  We study the classification of individuals defined by
a tuple $((x, x'), y)$, where $x\in \cX$ denotes a vector of
\emph{protected attributes}, $x'\in \cX'$ denotes a vector of
\emph{unprotected attributes}, and $y\in \{0, 1\}$ denotes a label. We
will write $X = (x, x')$ to denote the joint feature vector. We assume
that points $(X, y)$ are drawn i.i.d. from an unknown distribution
$\mathcal{P}$.  Let $D$ be a binary classifier, and let
$D(X) \in \{0,1\}$ denote the (possibly randomized) classification
induced by $D$ on individual $(X,y)$.

%When \emph{auditing} a fixed classifier $D$, it will be helpful to
%make reference to the distribution over examples $(X, y)$ together
%with their induced classification $D(X)$. Let $\PA(D)$ denote the
%induced \emph{target joint distribution} over the tuple $(x, y, D(X))$
% joint distribution} over the tuple $(x, x', y, D(X))$
%that results from sampling $(x, x', y) \sim \mathcal{P}$, and
%providing $x$, the true label $y$, and the classification
%$D(X) = D(x,x')$ but not the unprotected attributes $x'$.
%Note that
%the randomness here is over both the randomness of $\mathcal{P}$, and
%the potential randomness of the classifier $D$.


We will be concerned with learning and auditing classifiers $D$
satisfying a common statistical fairness constraint: equality
of false positive rates (also known as equal opportunity). The techniques in \citet{MSR} and \citet{KNRW18} also apply equally well to equality of false negative rates and equality of
classification rates (also known as statistical parity).\footnote{or more generally to any fairness constraint  that can be expressed as a linear equality on the conditional moments $\Ex{}{t(X, y, D(X) | \epsilon(X,y))},$ where $\epsilon(X,y)$ is an event defined with respect to $(X, y)$, and
$t: X \times \{0,1\} \times \{0,1\} \to [0,1]$ \cite{MSR}.  Equality of false positive rate is a particular instantiation of this kind of constraint where $\epsilon$  is the event $y = 0$, and $t = \textbf{1}\{D(X) =1\}$.}


Each fairness constraint is defined with
respect to a set of protected groups. We define sets of protected
groups via a family of indicator functions $\cF$ for those groups,
defined over protected attributes. Each $g:\cX \to \{0,1\} \in \cF$
has the semantics that $g(x) = 1$ indicates that an individual with
protected features $x$ is in group $g$.
%This notation generalizes the protected groups
%studied in \cite{MSR} which given a protected attribute $x'$ correspond to the set
%of functions $g_a = \textbf{1}\{x' = a\}$.
%Throughout the paper we refer to a classifier over $x, x'$ as
%\textit{marginally} fair, if it satisfies false positive strong subgroup fairness with respect to $\{g_a\}$.
 We now formally define
false positive subgroup fairness.


\begin{definition}[False Positive Subgroup Fairness]\label{fp-fair}
  Fix any classifier $D$, distribution $\mathcal{P}$, collection of
  group indicators $\cF$, and parameter $\gamma \in [0,1]$.  For each
  $g\in \cG$, define
\begin{align*}
  \fpsize(g, \cP) = \Pr_{\cP}[g(x) = 1, y=0],\\
  \fpdisp(g, D, \cP) = \left| \FP(D) - \FP(D, g) \right|
\end{align*}
  %\begin{align*}
  %  \fpsize(g, \cP) = \Pr_{\cP}[g(x) = 1, y=0] \quad \mbox{and,} \quad
  %                    \fpdisp(g, D, \cP) = \left| \FP(D) - \FP(D, g) \right|
  %\end{align*}
  where $\FP(D) = \Pr_{D, \cP} [D(X) = 1\mid y=0]$ and
  $\FP(D, g) = \Pr_{D, \cP}[D(X) = 1\mid g(x) =1 , y=0]$ denote the
  overall false-positive rate of $D$ and the false-positive rate of
  $D$ on group $g$ respectively.


  We say $D$
  satisfies $\gamma$-\emph{False Positive (FP) Fairness} with respect
  to $\mathcal{P}$ and $\cF$ if for every $g \in \cG$
  \[ \fpsize(g, \cP) \cdot \fpdisp(g, D, \cP) \leq \gamma.\] We will
  sometimes refer to $\FP(D)$ FP-base rate.% , which we write as:
  % $b_{FP} = b_{FP}(D, P) = \Pr[D(X) = 1 | y = 0].$
\end{definition}

Since we do not consider other measures in this paper, we refer to this notion as simply ``subgroup fairness.'' Given a fixed subgroup $g \in \cG$ we will refer to the quantity $\fpsize(g, \cP) \cdot \fpdisp(g, D, \cP)$ as the subgroup fairness wrt $g$, or alternately the $\gamma$-unfairness of $g$. The notion of subgroup fairness imposes a statistical constraint on combinatorially many groups definable by the protected attributes. This is in contrast to more common statistical fairness definitions, defined on coarse groups definable by a single protected attribute. Given a protected attribute $x_i$ and a value for that attribute $a$, define the function $g_{i,a}(x) = \textbf{1}\{x_i = a\}$ denoting the set of individuals who have that particular value of their protected attribute. In contrast to subgroup fairness, we refer to a classifier $D$ as
\textit{marginally} fair if it satisfies false positive subgroup fairness with respect to the functions $\{g_{i,a}\}$ for each protected attribute $x_i$ and realization $a$.

 If the
algorithm $D$ fails to satisfy the $\gamma$-subgroup fairness condition, then
we say that $D$ is $\gamma$-\emph{unfair} with respect to
$\mathcal{P}$ and $\cF$. We call any subgroup $g$ which witnesses this
unfairness a $\gamma$-\emph{unfair certificate} for
$(D, \mathcal{P})$.



An \emph{auditing algorithm} for a notion of fairness is given sample access to points from the underlying distribution, as well as the classification outcomes provided by $D$. It will either deem $D$ to be fair with respect to $\mathcal{P}$, or else produces a certificate of unfairness.
%
%
%\begin{definition}[Auditing Algorithm]
%  Fix a collection of group indicators $\cF$ over the
%  protected features, and any $\delta, \gamma, \gamma' \in (0, 1)$
%  such that $\gamma'\leq \gamma$.  A {$(\gamma, \gamma')$-auditing
%    algorithm for $\cF$ with respect to distribution $\mathcal{P}$} is
%  an algorithm $\cA$ such that for any classifier $D$, when given
%  access the distribution $\PA(D)$, $\cA$ runs in time
%  $\poly(1/\gamma', \log(1/\delta))$, and with probability $1-\delta$,
%  outputs a $\gamma'$-unfair certificate for $D$ whenever $D$ is
%  $\gamma$-unfair with respect to $\mathcal{P}$ and $\cF$. If $D$ is
%  $\gamma'$-fair, $\cA$ will output ``fair''.
%  % $D$ is $(\alpha, \beta - 2\rho)$-fair w.r.t. $P$ and $\cF$,
%  % then $A$ will output ``fair.''
%\end{definition}


The algorithms of \citet{MSR} and \citet{KNRW18} studied in this paper both assume access to oracles which can solve \emph{cost-sensitive classification
  (CSC)} problems. Formally, an instance of a CSC problem for the
class $\cH$ is given by a set of $n$ tuples
$\{(X_i, c_i^0, c_i^1)\}_{i=1}^n$ such that $c_i^\ell$ corresponds to
the cost for predicting label $\ell$ on point $X_i$. Given such an
instance as input, a CSC oracle finds a hypothesis $\hat h \in \cH$ that
minimizes the total cost across all points:
\begin{equation}
  \hat h \in \argmin_{h\in \cH} \sum_{i = 1}^n [h(X_i) c_i^1 + (1 -
  h(X_i)) c_i^0]\label{csc}
\end{equation}

Following both \cite{MSR} and \citet{KNRW18}, in all of the experiments in this paper we take the classes $\cH$ and $\cG$ to be linear threshold functions, and we use a linear regression heuristic for both auditing and learning. The heuristic finds a linear threshold function as follows: \\
\begin{itemize}
\item Train two linear regression models $r_0$, $r_1$ to predict $c_0$ and $c_1$ respectively.
\item Given a new point $x$, predict
the cost of classifying $x$ as $0$ and $1$ using our regression models: these are $r_0(x)$ and
$r_1(x)$ respectively.
\item Output the prediction $\hat{y}$ corresponding to lower predicted cost:
$\hat{y} = \argmin_{i \in \{0,1\}}r_i(x)$.
\end{itemize}
We leave the precise descriptions of the algorithm from \cite{KNRW18}
--- which we will refer to as the $\knrw$ algorithm --- to the
appendix. We refer the reader to \cite{KNRW18} for details about its
derivation and guarantees.\footnote{\cite{KNRW18}
  actually give two algorithms, one of which employs no-regret learning techniques and converges in a polynomial
  number of rounds, but is randomized; and the other of which
  is known to converge only in the limit (but is conjectured to converge
  quickly), and is deterministic. We focus on
  the deterministic algorithm in this paper, because it is more
  amenable to implementation, despite its weaker theoretical
  guarantees. We find that it performs well in practice despite its 
  weaker theory.}  At this point we
remark only that the algorithm operates by expressing the optimization
problem to be solved (minimize error, subject to subgroup fairness
constraints) as solving for the equilibrium in a two player zero-sum
game, between a \emph{Learner} and an \emph{Auditor}. The
\emph{Learner} has the set of hypothesis $\cH$ as its action (pure strategy) space,
and the \emph{Auditor} has the set of subgroups $\cF$ as its action
space. The best response problem for the Auditor corresponds to the
auditing problem: finding the subgroup $g \in \cF$ for which the
strategy of the learner violates the fairness constraints the
most. The best response problem for the Learner corresponds to solving
a weighted (but unconstrained) empirical risk minimization
problem. The best response problem for both players can be expressed
as solving a cost sensitive classification problem. The algorithm
$\knrw$ essentially simulates the \emph{fictitious play} of this game,
which proceeds over rounds, and in each round $t$ both players best
respond to their opponent's empirical history of play:
\begin{itemize} 
\item Learner plays $h_t$ in $\cH$ that minimizes objective function
  balancing error and unfairness on subgroups $g_1, \ldots, g_{t-1}$
  found by Auditor so far;
\item Auditor finds subgroup $g_t$ in $\cG$ on which the uniform
  distribution over $h_1, \ldots , h_t$ violates $\gamma$-fairness the
  most.
\end{itemize}
 This can be done efficiently assuming access to oracles which solve the cost sensitive classification problem over $\cF$ and $\cH$ respectively.

%%% Local Variables:
%%% mode: latex
%%% TeX-master: "main"
%%% End:
